% =========================================================================
\section{Puntos ordinarios: soluciones en serie entera}\label{sec:ordinarios}
% =========================================================================

% -------------------------------------------------------------------------
\subsection{Puntos ordinarios y puntos singulares}\label{sec:def-ordinario}
% -------------------------------------------------------------------------

\begin{definicion}[Punto ordinario y punto singular]\label{def:ordinario}
Un punto $x_0\in\R$ es un \emph{punto ordinario} de la ecuación
\eqref{eq:edo} si \emph{no anula} el polinomio $P$, es decir,
$P(x_0)\neq0$. En caso contrario, $P(x_0)=0$, se dice que $x_0$ es un
\emph{punto singular}.
\end{definicion}

\begin{ejemplo}\label{ej:ordinario}
En $(1-x^{2})\,y''-2x\,y'+\tfrac34\,y=0$ el punto $x_0=0$ es ordinario. De
hecho, todo $x_0\neq\pm1$ es ordinario, mientras que $x_0=\pm1$ son
singulares \cite{TEMAV}.
\end{ejemplo}

Si $x_0$ es ordinario podemos dividir por $P$ en un entorno de $x_0$.
Definimos
\begin{equation}\label{eq:pq}
  p(x)\defeq\frac{A(x)}{P(x)},\qquad
  q(x)\defeq\frac{B(x)}{P(x)},\qquad
  z\defeq x-x_0 ,
\end{equation}
con lo que la ecuación adopta la \emph{forma normal}
\begin{equation}\label{eq:normal}
  y''+p\,y'+q\,y=0 .
\end{equation}
Como $P(x_0)\neq0$, $p$ y $q$ son analíticas en $x_0$ por ser cocientes de
funciones analíticas con denominador no nulo (\cref{prop:algebra}):
\begin{equation}\label{eq:pq-serie}
  p=\sum_{k=0}^{\infty}p_k\,z^{k},\qquad
  q=\sum_{k=0}^{\infty}q_k\,z^{k},\qquad \abs{z}<R\defeq\min(R_p,R_q).
\end{equation}

\begin{observacion}[Una definición más general]\label{obs:ordinario-general}
Lo que interviene en toda la teoría es la analiticidad de $p$ y $q$ en
$x_0$, no el valor de $P(x_0)$. Por eso muchos textos llaman ordinario a
todo punto en el que $p$ y $q$ son analíticas \cite[§5.2]{BOYCE}. Por
ejemplo, en $x\,y''+(\sin x)\,y=0$ se tiene $P(0)=0$, pero
$q=\sin x/x$ tiene en $0$ una singularidad evitable
(\cref{ej:polos}\,(iv)): el punto $x_0=0$ es ordinario en sentido amplio y
todos los resultados de esta sección se le aplican.
\end{observacion}

% -------------------------------------------------------------------------
\subsection{El teorema de Fuchs}\label{sec:fuchs}
% -------------------------------------------------------------------------
El resultado central del tema afirma que en un punto ordinario el método de
series no sólo es formalmente viable, sino que produce soluciones
convergentes, y da una cota inferior explícita del radio
\cite[cap.~16]{RILEY}, \cite[cap.~5]{SIMMONS}.

\begin{teorema}[Fuchs, puntos ordinarios]\label{teo:fuchs}
Sea $x_0$ un punto ordinario de \eqref{eq:edo} y sean $p$, $q$ analíticas
en $\abs{z}<R$, con $z=x-x_0$ y $R=\min(R_p,R_q)$. Entonces:
\begin{enumerate}[label=\textup{(\roman*)}]
  \item para cada par $(a_0,a_1)\in\R^{2}$ existe una única serie
        $y=\sum_{n\ge0}a_n z^{n}$ que es solución de \eqref{eq:normal} con
        $y(x_0)=a_0$, $y'(x_0)=a_1$, y converge \emph{al menos} en
        $\abs{z}<R$;
  \item toda solución de \eqref{eq:normal} en $(x_0-R,\,x_0+R)$ es de esta
        forma; en particular, toda solución es analítica en $x_0$;
  \item si $A$ y $B$ son polinomios (o, más en general, funciones enteras),
        \begin{equation}\label{eq:cota-radio}
          R\;\ge\;\dist\bigl(x_0,\,\mathcal Z(P)\bigr),
          \qquad \mathcal Z(P)\defeq\{\zeta\in\C: P(\zeta)=0\},
        \end{equation}
        es decir, el radio es al menos la distancia, \emph{en el plano
        complejo}, de $x_0$ al cero de $P$ más cercano.
\end{enumerate}
\end{teorema}

La demostración se da en la \cref{sec:demo-fuchs}, una vez obtenida la
relación de recurrencia. La \cref{fig:fuchs} ilustra la cota
\eqref{eq:cota-radio}: es la versión, para las soluciones de la EDO, de lo
que el \cref{teo:radio} afirma para funciones racionales.

% Figura: disco de convergencia garantizado por el teorema de Fuchs
\begin{figure}[ht]
\centering
\begin{tikzpicture}[scale=1.15]
  % coordenadas: x0 = 0.6; ceros de P: zeta1 = 3.3 (real), zeta2,3 = -0.6 ± 1.9 i
  \def\xo{0.6}
  \pgfmathsetmacro{\radio}{sqrt((0.6+0.6)^2+1.9^2)}   % = dist(x0, zeta2) ≈ 2.247
  % disco
  \fill[cian, opacity=0.13] (\xo,0) circle (\radio);
  \draw[cian, line width=0.9pt] (\xo,0) circle (\radio);
  % ejes
  \draw[eje] (-2.6,0) -- (4.5,0) node[right, etiqgris] {$\Real z$};
  \draw[eje] (0,-2.75) -- (0,2.85) node[above, etiqgris] {$\Imag z$};
  % intervalo real de convergencia
  \draw[ambar, line width=2.4pt] ({\xo-\radio},0) -- ({\xo+\radio},0);
  \fill[ambar] ({\xo-\radio},0) circle (1.7pt);
  \fill[ambar] ({\xo+\radio},0) circle (1.7pt);
  \node[font=\footnotesize, text=ambar, anchor=north] at ({\xo-\radio-0.05},-0.12) {$x_0-R$};
  \node[font=\footnotesize, text=ambar, anchor=north] at ({\xo+\radio+0.05},-0.12) {$x_0+R$};
  % radio hasta el cero complejo más cercano
  \draw[blanco, line width=0.8pt, densely dashed] (\xo,0) -- (-0.6,1.9)
     node[midway, right=3pt, font=\small] {$R$};
  % centro
  \node[punto] at (\xo,0) {};
  \node[font=\small, anchor=north] at (\xo+0.12,-0.12) {$x_0$};
  % ceros de P
  \node[polo] at (3.3,0) {};
  \node[polo] at (-0.6,1.9) {};
  \node[polo] at (-0.6,-1.9) {};
  \node[font=\small, text=rojo, anchor=south] at (3.3,0.14) {$\zeta_1$};
  \node[font=\small, text=rojo, anchor=east] at (-0.72,1.98) {$\zeta_2$};
  \node[font=\small, text=rojo, anchor=east] at (-0.72,-1.98) {$\zeta_3=\overline{\zeta_2}$};
  % rótulo del disco
  \node[font=\small, text=cian, anchor=west] at (2.15,1.75) {$\abs{z-x_0}<R$};
  % leyenda
  \begin{scope}[shift={(5.0,1.6)}]
    \node[polo] at (0,0) {};
    \node[etiqgris, anchor=west] at (0.2,0) {ceros de $P$ en $\C$};
    \draw[ambar, line width=2.4pt] (-0.15,-0.55) -- (0.15,-0.55);
    \node[etiqgris, anchor=west] at (0.2,-0.55) {convergencia en $\R$};
  \end{scope}
\end{tikzpicture}
\caption{Teorema de Fuchs. El disco de convergencia garantizado está
centrado en el punto ordinario $x_0$ y llega hasta el cero complejo de $P$
más próximo, aquí el par conjugado $\zeta_{2,3}$. Aunque el único cero
\emph{real}, $\zeta_1$, está más lejos, es $\zeta_2$ quien fija
$R=\abs{\zeta_2-x_0}$ y, con él, el intervalo real $(x_0-R,\,x_0+R)$ en
el que converge la serie.}
\label{fig:fuchs}
\end{figure}


\begin{observacion}[El radio garantizado es una cota inferior]\label{obs:cota}
La cota \eqref{eq:cota-radio} puede no alcanzarse: si una de las series
solución se trunca en un polinomio, ésta converge en todo $\R$ aunque $P$
tenga ceros (polinomios de Legendre, \cref{sec:legendre}). Lo que el teorema
asegura es que \emph{ninguna} solución deja de converger antes de llegar al
cero complejo de $P$ más próximo.
\end{observacion}

% -------------------------------------------------------------------------
\subsection{Ansatz y derivación término a término}\label{sec:ansatz}
% -------------------------------------------------------------------------
Proponemos
\begin{equation}\label{eq:ansatz}
  y=\sum_{n=0}^{\infty}a_n\,z^{n}.
\end{equation}
Derivando término a término (\cref{teo:series}),
\begin{align}
  y'  &=\sum_{n=1}^{\infty}n\,a_n\,z^{n-1}, \label{eq:yp}\\
  y'' &=\sum_{n=2}^{\infty}n(n-1)\,a_n\,z^{n-2}. \label{eq:ypp}
\end{align}
Los límites inferiores suben porque el término $n=0$ de $y'$ (y los
$n=0,1$ de $y''$) son nulos. De momento el cálculo es \emph{formal}: la
convergencia de \eqref{eq:ansatz} se establecerá en la \cref{sec:demo-fuchs},
y es entonces cuando el \cref{teo:series} legitima a posteriori cada paso.

% -------------------------------------------------------------------------
\subsection{Sustitución y cambio de índices}\label{sec:indices}
% -------------------------------------------------------------------------
Introduciendo \eqref{eq:pq-serie}, \eqref{eq:yp} y \eqref{eq:ypp} en
\eqref{eq:normal}:
\begin{equation}\label{eq:sustitucion}
  \sum_{n=2}^{\infty}n(n-1)\,a_n z^{n-2}
  +\Bigl(\sum_{k=0}^{\infty}p_k z^{k}\Bigr)\Bigl(\sum_{n=1}^{\infty}n\,a_n z^{n-1}\Bigr)
  +\Bigl(\sum_{k=0}^{\infty}q_k z^{k}\Bigr)\Bigl(\sum_{n=0}^{\infty}a_n z^{n}\Bigr)=0 .
\end{equation}
Queremos escribir todas las series en potencias $z^{n}$. El índice de
suma es \emph{mudo} y puede desplazarse libremente:
\begin{equation}\label{eq:desplazamiento}
  \sum_{n=m}^{\infty}c_n\,z^{n-m}
  \;\overset{n\,\mapsto\,n+m}{=}\;
  \sum_{n=0}^{\infty}c_{n+m}\,z^{n}.
\end{equation}
\begin{itemize}
  \item En $y''$ tomamos $m=2$ (es decir, llamamos $n$ a $n-2$):
        \begin{equation}\label{eq:ypp-desplazada}
          y''=\sum_{n=0}^{\infty}(n+2)(n+1)\,a_{n+2}\,z^{n}.
        \end{equation}
  \item En $y'$ tomamos $m=1$:
        \begin{equation}\label{eq:yp-desplazada}
          y'=\sum_{n=0}^{\infty}(n+1)\,a_{n+1}\,z^{n}.
        \end{equation}
\end{itemize}
Ahora las tres series empiezan en $n=0$ y llevan $z^{n}$; usando un índice
$n'$ en los factores para no confundirlo con el de los productos,
\begin{equation}\label{eq:sustitucion2}
  \sum_{n}(n+2)(n+1)a_{n+2}z^{n}
  +\Bigl(\sum_{k}p_k z^{k}\Bigr)\Bigl(\sum_{n'}(n'+1)a_{n'+1}z^{n'}\Bigr)
  +\Bigl(\sum_{k}q_k z^{k}\Bigr)\Bigl(\sum_{n'}a_{n'}z^{n'}\Bigr)=0 .
\end{equation}

% -------------------------------------------------------------------------
\subsection{Producto de Cauchy}\label{sec:cauchy}
% -------------------------------------------------------------------------

\begin{lema}[Producto de Cauchy]\label{lem:cauchy}
Si las series $\sum_{k\ge0}c_kz^{k}$ y $\sum_{n'\ge0}b_{n'}z^{n'}$ convergen
absolutamente en $\abs{z}<R$, entonces, en ese disco,
\begin{equation}\label{eq:cauchy}
  \Bigl(\sum_{k=0}^{\infty}c_k z^{k}\Bigr)\Bigl(\sum_{n'=0}^{\infty}b_{n'}z^{n'}\Bigr)
  =\sum_{n=0}^{\infty}\Bigl(\sum_{k=0}^{n}c_k\,b_{n-k}\Bigr)z^{n},
\end{equation}
y la serie del miembro derecho converge también absolutamente.
\end{lema}

\begin{proof}
Al multiplicar aparecen los términos $c_kb_{n'}z^{k}z^{n'}=c_kb_{n'}z^{k+n'}$
para todos los pares $(k,n')\in\N^{2}$. La serie doble correspondiente
converge absolutamente, pues
$\sum_{k,n'}\abs{c_k}\abs{b_{n'}}\abs{z}^{k+n'}
=\bigl(\sum_k\abs{c_k}\abs{z}^k\bigr)\bigl(\sum_{n'}\abs{b_{n'}}\abs{z}^{n'}\bigr)<\infty$,
y por tanto puede sumarse en cualquier orden. Agrupamos los términos con la
misma potencia total $n=k+n'$, es decir, recorremos las diagonales de la
\cref{fig:cauchy}. Fijado $n$, se tiene $n'=n-k$, y como $n'\ge0$, el
índice $k$ sólo recorre $0,1,\dots,n$.
\end{proof}

% Figura: producto de Cauchy como suma por diagonales k + n' = n
\begin{figure}[ht]
\centering
\begin{tikzpicture}[x=0.82cm, y=0.82cm]
  % diagonales k + n' = n
  \foreach \n in {1,...,6}{
    \draw[guia] (\n+0.3,-0.3) -- (-0.3,\n+0.3);}
  % diagonal destacada n = 4
  \draw[cian, line width=1.2pt] (4.35,-0.35) -- (-0.35,4.35);
  % ejes
  \draw[eje] (-0.4,0) -- (6.6,0) node[right, etiqgris] {$k$};
  \draw[eje] (0,-0.4) -- (0,6.6) node[above, etiqgris] {$n'$};
  \foreach \t in {0,...,5}{
    \node[etiqgris, anchor=north] at (\t,-0.45) {$\t$};
    \node[etiqgris, anchor=east] at (-0.45,\t) {$\t$};}
  % malla de pares (k, n')
  \foreach \k in {0,...,5}{\foreach \m in {0,...,5}{
    \fill[grismedio] (\k,\m) circle (2.2pt);}}
  \foreach \k in {0,...,4}{
    \pgfmathtruncatemacro{\m}{4-\k}
    \fill[cian] (\k,\m) circle (3.4pt);}
  % rótulo de la diagonal
  \node[nodocian, anchor=west, align=left] (r) at (7.4,4.1)
    {\textcolor{cian}{\bfseries Diagonal $n=4$}\\[2pt]
     $c_0b_4+c_1b_3+c_2b_2+c_3b_1+c_4b_0$\\[1pt]
     \footnotesize\color{gris} coeficiente de $z^{4}$ en el producto};
  \draw[flecha, cian] (r.west) to[out=180, in=45] (2.15,2.15);
  \node[nodo, anchor=west, align=left] at (7.4,1.4)
    {\footnotesize\color{gris} Cada punto $(k,n')$ es el término\\
     \footnotesize\color{gris} $c_k b_{n'}z^{k+n'}$; en la diagonal\\
     \footnotesize\color{gris} $k+n'=n$ todos llevan $z^{n}$.};
\end{tikzpicture}
\caption{Producto de Cauchy de dos series de potencias. Los términos
$c_kb_{n'}z^{k+n'}$ se agrupan por diagonales $k+n'=n$; como $n'=n-k\ge0$,
en la diagonal $n$ el índice $k$ sólo recorre $0,1,\dots,n$, lo que da el
coeficiente $\sum_{k=0}^{n}c_kb_{n-k}$ de \eqref{eq:cauchy}.}
\label{fig:cauchy}
\end{figure}


Aplicando el \cref{lem:cauchy} con $b_{n'}=(n'+1)\,a_{n'+1}$, de modo que
$b_{n-k}=(n-k+1)\,a_{n-k+1}$, y con $b_{n'}=a_{n'}$:
\begin{align}
  p\,y' &=\sum_{n=0}^{\infty}\Bigl(\sum_{k=0}^{n}(n-k+1)\,p_k\,a_{n-k+1}\Bigr)z^{n},\label{eq:pyp}\\
  q\,y  &=\sum_{n=0}^{\infty}\Bigl(\sum_{k=0}^{n}q_k\,a_{n-k}\Bigr)z^{n}.\label{eq:qy}
\end{align}

% -------------------------------------------------------------------------
\subsection{Identificación de coeficientes: la relación de recurrencia}\label{sec:recurrencia}
% -------------------------------------------------------------------------
Toda la ecuación es ahora una única serie $\sum_n C_n z^{n}=0$ válida para
todo $\abs{z}<R$. Por la unicidad del desarrollo en serie de potencias
(\cref{cor:unicidad}) ---la serie nula sólo tiene coeficientes nulos---,
cada coeficiente debe anularse.

\begin{resultado}[Relación de recurrencia en un punto ordinario]\label{res:recurrencia}
Los coeficientes de toda solución \eqref{eq:ansatz} de \eqref{eq:normal}
satisfacen, para todo $n\ge0$,
\begin{equation}\label{eq:recurrencia}
  (n+2)(n+1)\,a_{n+2}
  +\sum_{k=0}^{n}\Bigl[(n-k+1)\,p_k\,a_{n-k+1}+q_k\,a_{n-k}\Bigr]=0 .
\end{equation}
\end{resultado}

Ésta es la ecuación en diferencias anunciada en la \cref{sec:por-que}: la
EDO se ha convertido en una relación de recurrencia para los coeficientes.
Si $p$ y $q$ son polinomios de grado $\le d$, la suma en \eqref{eq:recurrencia}
tiene a lo sumo $d+1$ términos no nulos y la recurrencia es de orden finito;
en general involucra todos los coeficientes anteriores.

\begin{observacion}[Erratas frecuentes]\label{obs:erratas}
Al transcribir \eqref{eq:recurrencia} conviene vigilar dos detalles en los
que es fácil equivocarse: en el término de $q$ la potencia que acompaña a
$q_k$ es $z^{k}$, no $z^{2}$; y en el de $p$ el subíndice correcto es
$a_{n-k+1}$, multiplicado por el factor $(n-k+1)$ procedente de la derivada.
\end{observacion}

% -------------------------------------------------------------------------
\subsection{Resolución de la recurrencia}\label{sec:resolucion}
% -------------------------------------------------------------------------
Como $(n+2)(n+1)\neq0$ para $n\ge0$, podemos despejar el coeficiente de
mayor índice:
\begin{equation}\label{eq:despeje-an2}
  \boxed{\;a_{n+2}=-\frac{1}{(n+2)(n+1)}
  \sum_{k=0}^{n}\Bigl[(n-k+1)\,p_k\,a_{n-k+1}+q_k\,a_{n-k}\Bigr],
  \qquad n\ge0.\;}
\end{equation}
El miembro derecho sólo contiene $a_0,a_1,\dots,a_{n+1}$: conocidos todos los
anteriores, $a_{n+2}$ queda determinado (\cref{fig:dependencia}).

\begin{proposicion}[Determinación por inducción]\label{prop:induccion}
Para cada par $(a_0,a_1)$ existe una única sucesión $(a_n)_{n\ge0}$ que
cumple \eqref{eq:despeje-an2}. Además, cada $a_n$ es una combinación lineal
de $a_0$ y $a_1$.
\end{proposicion}

\begin{proof}
Por inducción fuerte sobre $n$: \eqref{eq:despeje-an2} define $a_{n+2}$ a
partir de $a_0,\dots,a_{n+1}$, ya determinados. La linealidad se debe a que
\eqref{eq:despeje-an2} es lineal en los $a_j$.
\end{proof}

% Figura: estructura de dependencia de la recurrencia general
\begin{figure}[ht]
\centering
\begin{tikzpicture}
  \tikzset{coef/.style={circle, draw=grismedio, fill=cajafondo, line width=0.6pt,
                        minimum size=0.95cm, inner sep=0pt, font=\small}}
  \foreach \i in {0,...,6}{
    \node[coef] (a\i) at (1.75*\i,0) {$a_{\i}$};}
  \node[font=\small, text=gris] at (1.75*7-0.35,0) {$\cdots$};
  % datos libres
  \node[coef, draw=ambar, line width=1pt] at (a0) {$a_{0}$};
  \node[coef, draw=ambar, line width=1pt] at (a1) {$a_{1}$};
  % coeficiente calculado
  \node[coef, draw=cian, line width=1.2pt] at (a4) {$a_{4}$};
  % flechas de dependencia hacia a4 (n = 2)
  \foreach \j/\h in {0/62, 1/52, 2/40, 3/30}{
    \draw[flecha, cian!75] (a\j.north) to[out=\h, in=180-\h+8] (a4.north);}
  % llave bajo los datos libres
  \draw[decorate, decoration={brace, amplitude=5pt, mirror}, ambar, line width=0.6pt]
    ($(a0.south west)+(0,-0.15)$) -- ($(a1.south east)+(0,-0.15)$)
    node[midway, below=7pt, font=\footnotesize, text=ambar, align=center]
    {datos libres\\$a_0=y(x_0)$,\ $a_1=y'(x_0)$};
  \draw[decorate, decoration={brace, amplitude=5pt, mirror}, grismedio, line width=0.6pt]
    ($(a2.south west)+(0,-0.15)$) -- ($(a6.south east)+(0,-0.15)$)
    node[midway, below=7pt, font=\footnotesize, text=gris, align=center]
    {determinados por \eqref{eq:despeje-an2},\\uno tras otro};
  \node[font=\footnotesize, text=cian, anchor=west, align=left] at ($(a5.north)+(-0.1,0.95)$)
    {$a_{4}$ depende de $a_0,a_1,a_2,a_3$\\($n=2$ en la recurrencia)};
\end{tikzpicture}
\caption{Estructura de la recurrencia \eqref{eq:despeje-an2}: el coeficiente
$a_{n+2}$ depende de \emph{todos} los anteriores (aquí, $a_4$). Sólo $a_0$ y
$a_1$ quedan libres; son las dos constantes de integración.}
\label{fig:dependencia}
\end{figure}


\paragraph{Primeros términos}
Particularizando \eqref{eq:despeje-an2}:
\begin{itemize}
  \item $n=0$ (sólo $k=0$):
        \begin{equation}\label{eq:a2}
          a_2=-\tfrac12\bigl(p_0a_1+q_0a_0\bigr).
        \end{equation}
  \item $n=1$ ($k=0,1$):
        \begin{equation}\label{eq:a3}
          a_3=-\tfrac16\bigl(2p_0a_2+q_0a_1+p_1a_1+q_1a_0\bigr).
        \end{equation}
  \item $n=2$ ($k=0,1,2$):
        \begin{equation}\label{eq:a4}
          a_4=-\tfrac1{12}\bigl(3p_0a_3+(q_0+2p_1)\,a_2+(q_1+p_2)\,a_1+q_2a_0\bigr).
        \end{equation}
\end{itemize}

\begin{observacion}[Comprobación]\label{obs:comprobacion}
La fórmula \eqref{eq:a2} puede verificarse sin series: evaluando
\eqref{eq:normal} en $x_0$ se obtiene
$y''(x_0)=-p(x_0)\,y'(x_0)-q(x_0)\,y(x_0)=-p_0a_1-q_0a_0$, y por el
\cref{teo:series}\,(iii), $a_2=y''(x_0)/2$. Del mismo modo, derivando la
ecuación sucesivamente se recupera cada $a_n=y^{(n)}(x_0)/n!$: la
recurrencia es el algoritmo de Taylor organizado.
\end{observacion}

% -------------------------------------------------------------------------
\subsection{Condiciones iniciales y sistema fundamental}\label{sec:fundamental}
% -------------------------------------------------------------------------
Evaluando la serie \eqref{eq:ansatz} y su derivada \eqref{eq:yp} en $z=0$,
\begin{equation}\label{eq:a0a1}
  a_0=y(x_0),\qquad a_1=y'(x_0).
\end{equation}
Son las condiciones iniciales: las dos constantes de integración que se
esperan en una EDO de segundo orden. La recurrencia no las fija; quedan
libres.

Como la ecuación es lineal, basta elegir dos juegos de condiciones
iniciales:
\begin{equation}\label{eq:y1y2}
  y_1\defeq y\big|_{(a_0,a_1)=(1,0)},\qquad
  y_2\defeq y\big|_{(a_0,a_1)=(0,1)}.
\end{equation}

\begin{proposicion}[Sistema fundamental]\label{prop:fundamental}
Las funciones $y_1$, $y_2$ forman un sistema fundamental de soluciones de
\eqref{eq:normal} en $(x_0-R,\,x_0+R)$, y la solución general es
\begin{equation}\label{eq:general}
  y=a_0\,y_1+a_1\,y_2,\qquad a_0,a_1\in\R .
\end{equation}
\end{proposicion}

\begin{proof}
Su wronskiano en $x_0$ es
\begin{equation}\label{eq:wronskiano}
  W(x_0)=\begin{vmatrix} y_1(x_0) & y_2(x_0)\\ y_1'(x_0) & y_2'(x_0)\end{vmatrix}
  =\begin{vmatrix}1&0\\0&1\end{vmatrix}=1\neq0 .
\end{equation}
Además $W'=y_1y_2''-y_1''y_2=-p\,W$, de donde (fórmula de Abel)
$W(x)=\exp\bigl(-\int_{x_0}^{x}p(t)\,\dif t\bigr)\neq0$ en todo el
intervalo: $y_1$ e $y_2$ son linealmente independientes. Por la
\cref{prop:induccion}, la solución con datos $(a_0,a_1)$ tiene coeficientes
lineales en $(a_0,a_1)$, luego es $a_0y_1+a_1y_2$.
\end{proof}

% -------------------------------------------------------------------------
\subsection{Demostración del teorema de Fuchs}\label{sec:demo-fuchs}
% -------------------------------------------------------------------------
La única afirmación no trivial del \cref{teo:fuchs} es que la serie formal
construida con \eqref{eq:despeje-an2} \emph{converge}. Lo probaremos con el
\emph{método de las series mayorantes} de Cauchy: acotaremos $\abs{a_n}$ por
una sucesión $b_n$ cuyo comportamiento conocemos exactamente
\cite[cap.~3]{CODDINGTON}, \cite[§4.1]{TESCHL}.

\begin{proof}[Demostración del \cref{teo:fuchs}]
Si $a_0=a_1=0$, todos los $a_n$ son nulos y no hay nada que probar;
supongamos $(a_0,a_1)\neq(0,0)$.

\emph{Paso 1: cotas de Cauchy.} Fijemos $0<r<R$. Como $\sum p_kr^{k}$ y
$\sum q_kr^{k}$ convergen, sus términos están acotados: existe $M>0$ tal que
\begin{equation}\label{eq:cotas-cauchy}
  \abs{p_k}\le M\,r^{-k},\qquad \abs{q_k}\le M\,r^{-k}\qquad\forall k\ge0 .
\end{equation}

\emph{Paso 2: sucesión mayorante.} Definimos $b_0=\abs{a_0}$,
$b_1=\abs{a_1}$ y, para $n\ge0$,
\begin{equation}\label{eq:mayorante}
  (n+2)(n+1)\,b_{n+2}=M\sum_{j=0}^{n}r^{\,j-n}\bigl[(j+1)\,b_{j+1}+b_j\bigr]+M\,r\,b_{n+1}.
\end{equation}
Veamos por inducción que $\abs{a_n}\le b_n$. Tomando módulos en
\eqref{eq:recurrencia}, usando \eqref{eq:cotas-cauchy} y cambiando
$j=n-k$,
\begin{align*}
  (n+2)(n+1)\abs{a_{n+2}}
  &\le\sum_{k=0}^{n}\bigl[(n-k+1)\abs{p_k}\abs{a_{n-k+1}}+\abs{q_k}\abs{a_{n-k}}\bigr]\\
  &\le M\sum_{j=0}^{n}r^{\,j-n}\bigl[(j+1)\,b_{j+1}+b_j\bigr]
  \;\le\;(n+2)(n+1)\,b_{n+2}.
\end{align*}

\emph{Paso 3: el cociente $b_{n+2}/b_{n+1}$.} Escribiendo
\eqref{eq:mayorante} para $n$ (multiplicada por $r$) y para $n-1$, con
$n\ge1$, y restando, todos los términos de la suma salvo el de $j=n$ se
cancelan:
\[
  r\,(n+2)(n+1)\,b_{n+2}-(n+1)n\,b_{n+1}
  =M r\bigl[(n+1)b_{n+1}+b_n\bigr]+Mr^{2}b_{n+1}-Mr\,b_n ,
\]
es decir,
\begin{equation}\label{eq:cociente}
  \frac{b_{n+2}}{b_{n+1}}
  =\frac{n(n+1)+Mr\,(n+1)+Mr^{2}}{r\,(n+2)(n+1)}
  \;\xrightarrow[n\to\infty]{}\;\frac1r .
\end{equation}
(Todos los $b_n$ con $n\ge2$ son positivos: $b_2=\tfrac M2(b_0+b_1+rb_1)>0$
y \eqref{eq:cociente} propaga la positividad.)

\emph{Paso 4: convergencia.} Por el criterio del cociente,
$\sum b_n\abs{z}^{n}$ converge si $\abs{z}<r$, y por comparación
$\sum a_nz^{n}$ converge absolutamente en $\abs{z}<r$. Como $r<R$ es
arbitrario, la serie converge en $\abs{z}<R$.

\emph{Paso 5: es solución.} En $\abs{z}<R$ son lícitas la derivación término a
término (\cref{teo:series}) y el producto de Cauchy (\cref{lem:cauchy}); el
coeficiente de $z^{n}$ en $y''+py'+qy$ es exactamente el miembro izquierdo de
\eqref{eq:recurrencia}, que es nulo. Además $y(x_0)=a_0$ e $y'(x_0)=a_1$.
Esto prueba (i).

\emph{Paso 6: toda solución es una serie.} Si $\varphi$ es una solución en
$(x_0-R,x_0+R)$, la serie $y$ con $a_0=\varphi(x_0)$, $a_1=\varphi'(x_0)$ es
otra solución con los mismos datos, y por el \cref{teo:eyu}, $\varphi=y$.
Esto prueba (ii).

\emph{Paso 7: cota mediante los ceros de $P$.} Si $A$ y $B$ son enteras,
$p=A/P$ y $q=B/P$ son holomorfas en el disco
$\disco{x_0}{d}$, $d=\dist(x_0,\mathcal Z(P))$, y por el \cref{teo:radio}
sus series de Taylor convergen en él: $R_p,R_q\ge d$. Esto prueba (iii).
\end{proof}
